\section*{Conclusion Personnelle}

\subsection*{Conclusion Personnelle - Romain}

Ce projet de protection des données m’a offert une expérience enrichissante à plusieurs niveaux.

Tout d'abord, le fait de travaille en groupe m’a permis de développer mes compétences collaboratives, notamment sur la manière de se répartir le travail et d'utiliser les outils pour collaborer tels que Github de manière efficace.

De plus, j'ai pu approfondir mes connaissances en science des données et mieux comprendre l'intérêt de pré-traiter ses données pour un projet similaire.

Enfin, ce projet m'a permis de comprendre l'intérêt de vulgariser les résultats de manière claire et accessible. Par ailleurs, j'ai aussi découvert l'outil streamlit qui peut s'avérer très utile pour présenter des résultats de manière interactive.

\subsection*{Conclusion Personnelle - Jean-Raphael}

Le projet fut enrichissant de connaissances, j'ai découvert l'utilisation de streamlit et j'ai pu apprendre à fouiller dans un dataset pour tirer des ressources intérressantes.

De mon côté, je me suis occupé de la partie d'analyse des données du côté network, c'est-à-dire que j'ai fait la réalisation de tous les plots d'analyses pour le projet.

J'ai également réalisé l'implémentation de CART,KNN,random forest et Xgboost pour le côté network ce qui m'a permis de manipuler les paramètres de ces algorithmes.

Ce projet m'a apporté des solides connaissances dans l'analyse de données via des arbres de décisions et permet également de se questionner sur les métriques à utiliser pour évaluer ses modèles.

De plus, ce projet m'a permis de mettre en lumière l'importance de choisir ses données correctement et de comprendre pourquoi le modèle agit d'une telle façon.

\subsection*{Conclusion Personnelle - Bergamini Enzo}

Ce projet a été une opportunité d’explorer et de comparer divers algorithmes de classification supervisée, tels que KNN, CART, Random Forest et XGBoost, dans le cadre d’un problème impliquant des classes majoritaires et minoritaires. 

Les résultats obtenus ont mis en lumière la capacité de chaque modèle à gérer des jeux de données déséquilibrés, tout en révélant leurs atouts et leurs limites respectives. Ce travail m’a permis de renforcer mes compétences dans la manipulation de données, l’évaluation de modèles et l’interprétation des métriques de classification. 

J’ai également compris l’importance de choisir un algorithme en prenant en compte non seulement ses performances globales, mais aussi des contraintes spécifiques à l’application, comme le temps de calcul, la consommation mémoire. Au cours de ce projet, ma contribution s’est concentrée sur l’entraînement des algorithmes de machine learning, l’interprétation de leurs résultats et la validation de leurs performances.

\subsection*{Conclusion Personnelle - Verdon Johan}

Ce projet a été stimulant et enrichissant pour moi. Dès les premières étapes, j’ai pris l’initiative de développer les fichiers nécessaires au nettoyage des données et de mener une première analyse exploratoire pour mieux comprendre les informations à traiter. Progressivement, j’ai structuré le projet, en lançant également le dépôt Git sur lequel l’ensemble de l’équipe a collaboré. Cette anticipation m’a permis de faciliter l’organisation générale : chacun a pu rapidement trouver sa place et s’assigner des tâches utiles au projet. J’ai pris plaisir à assumer ce rôle informel de coordination et de gestion du projet dans son ensemble.

Sur le plan technique, j’ai contribué de manière significative en créant les scripts de nettoyage des données physiques, en menant une analyse approfondie et en développant les modèles de classification. De plus, j’ai conçu une application Streamlit permettant de visualiser et d’exploiter ces résultats pour les données physiques. Parallèlement, j’ai aidé mes coéquipiers sur les données réseaux, la structuration du dépôt Git et la rédaction du rapport final.

J’ai particulièrement apprécié les défis posés par les données et les objectifs du projet. Ce rôle global que j’ai endossé, mêlant à la fois aspects techniques et organisationnels, m’a permis de m’investir pleinement tout en développant des compétences en gestion de projet et en collaboration. Je suis satisfait de la manière dont nous avons progressé collectivement pour mener à bien cette mission.
